% !TEX encoding = UTF-8 Unicode

\providecommand{\cleartorecto}{\cleardoublepage}
\providecommand{\adjustmtc}{}

\providecommand{\Frangi}{\textnormal{\textsc{Frangi}}}
\providecommand{\Hess}{\mathcal H}
\providecommand{\matnorm}[1]{\left\lvert\mkern-2mu\left\lvert\mkern-2mu\left\lvert #1 \right\rvert\mkern-2mu\right\rvert\mkern-2mu\right\rvert}

\providecommand{\Lab}{\mathsf S}
\providecommand{\Conf}{\mathcal S}
\providecommand{\Obs}{\mathcal O}
\providecommand{\Wspace}{\mathsf W}
\providecommand{\RG}{\operatorname{RG}}
\providecommand{\ov}{\operatorname{ov}}

%\providecommand{\HGPClusterer}{\textsc{HGP-Clusterer}}
\providecommand{\DBSCAN}{\textsc{DBSCAN}}
\providecommand{\HDBSCAN}{\textsc{HDBSCAN}}
\providecommand{\KNN}{$K$-\textsc{NN}}

\makeatletter
\providecommand{\Needspace}[1]{%
  \par
  \@tempskipa=#1\relax
  \ifdim\pagegoal=\maxdimen
  \else
    \ifdim\dimexpr\pagegoal-\pagetotal-\pageshrink\relax<\@tempskipa
      \newpage
    \fi
  \fi
}
\makeatother

\cleartorecto
\Chapter{Notations mathématiques}
\label{chap:notations}

\begingroup
\small
\setlength{\tabcolsep}{4pt}
\renewcommand{\arraystretch}{1.17}

\newcommand{\ChapRef}[1]{\hyperref[#1]{chap.~\ref*{#1}}}
\newcommand{\ChapRefs}[2]{\ChapRef{#1}, \ChapRef{#2}}
\newcommand{\ChapRefsThree}[3]{\ChapRef{#1}, \ChapRef{#2}, \ChapRef{#3}}
\newcommand{\CtxIntro}{\ChapRef{sec:Introduction}}
\newcommand{\CtxGlobal}{global}

\newenvironment{NotationTable}[2][0.30\textheight]{%
  \Needspace{#1}
  \begin{tcolorbox}[
    enhanced,
    unbreakable,
    sharp corners,
    boxrule=0.45pt,
    colback=white,
    colframe=black!32,
    colbacktitle=black!7,
    coltitle=black,
    fonttitle=\bfseries,
    title={#2},
    before skip=8pt,
    after skip=8pt,
    left=1.4mm,
    right=1.4mm,
    top=1mm,
    bottom=1mm
  ]
  \begin{tabular}{@{}>{\raggedleft\arraybackslash}p{0.225\textwidth}
                  p{0.565\textwidth}
                  >{\raggedright\arraybackslash}p{0.145\textwidth}@{}}
  \toprule
  \rowcolor{black!6}
  \textbf{Symbole} & \textbf{Signification} & \textbf{Chapitre(s)}\\
  \midrule
}{%
  \bottomrule
  \end{tabular}
  \end{tcolorbox}
}

\newcommand{\Not}[3]{%
  \(\displaystyle #1\) & #2 & {\footnotesize #3}\\
}

\newcommand{\NotText}[3]{%
  #1 & #2 & {\footnotesize #3}\\
}

\begin{NotationTable}{Abréviations}
\Not{\Fig}{Figure.}{\CtxGlobal}
\Not{\Tab}{Tableau.}{\CtxGlobal}
\Not{\Def}{Définition (avec référence dans le cadre si la définition n'est pas originale).}{\CtxGlobal}
\Not{\Res}{Fait déjà connu.}{\CtxGlobal}
\Not{\Fact \text{ et }\Th}{Résultats originaux.}{\CtxGlobal}
\Not{\S}{Section.}{\CtxGlobal}
\end{NotationTable}

\begin{NotationTable}{Conventions générales}
\Not{\eqdef}{égal à (par définition).}{\CtxGlobal}
\Not{\mathbb N^\star}{ensemble des entiers naturels non nuls.}{\CtxGlobal}
\Not{\mathbb R^p}{espace euclidien de dimension \(p\).}{\CtxGlobal}
\Not{p}{dimension ambiante.}{\CtxGlobal}
\Not{n}{nombre de points ou de sommets, selon le contexte.}{\CtxGlobal}
\Not{x,x',y,z}{points, pixels ou sommets selon le contexte ; en général \(x,x'\in\mathcal X\) et \(y\in\mathbb R^p\).}{\CtxGlobal}
\Not{d(x,x')}{distance entre deux points.}{\CtxGlobal}
\Not{\|x-x'\|}{distance euclidienne.}{\CtxGlobal}
\Not{\overline B(x,r)}{boule fermée de centre \(x\) et de rayon \(r\).}{\CtxGlobal}
%\Not{B(x,r)}{boule lorsque l'ouverture ou la fermeture ne change pas l'argument.}{\CtxGlobal}
\Not{\partial A,\ \mathring A}{frontière et intérieur de l'ensemble \(A\).}{\CtxGlobal}
\end{NotationTable}

\begin{NotationTable}{Ensembles, mesure et connexité}
\Not{|A|}{cardinal de \(A\), ou mesure de \nom{Lebesgue} de $A$ selon le contexte.}{\CtxGlobal}%si \(A\subset\mathbb R^p\) est mesurable
\Not{\omega_p}{volume de la boule unité de \(\mathbb R^p\).}{\ChapRef{chap:single-linkage_trois_vues}}
\Not{\mathds 1_{\{E\}}}{indicatrice de l'événement \(E\).}{\CtxGlobal}
\Not{\mathbb P,\ \mathbb E}{probabilité et espérance.}{\CtxGlobal}
\Not{\mathcal O(\cdot)}{ordre de grandeur asymptotique.}{\CtxGlobal}
\Not{\pi_0(A)}{ensemble des composantes connexes de \(A\).}{\ChapRefs{chap:single-linkage_trois_vues}{chap:hgp_clusterer}}
%\Not{\mathrm{cc}}{type de composantes connexes considéré.}{\ChapRef{sec:chap_percolation}}
\Not{\delta_r(A)}{dilatation de \(A\) par une boule de rayon \(r\).}{\ChapRefs{chap:single-linkage_trois_vues}{sec:chap_percolation}}
\Not{A\dot\cup B}{union disjointe.}{\ChapRef{chap:axiomatique}}
\end{NotationTable}

\begin{NotationTable}{Partitions et hiérarchies}
\Not{\mathcal X}{nuage fini de points.}{\CtxIntro}
\Not{P}{partition d'un ensemble.}{\ChapRef{chap:axiomatique}}
\Not{\Partitions{\mathcal X}}{ensemble des partitions de \(\mathcal X\).}{\ChapRef{chap:axiomatique}}
\Not{\sim}{relation d'équivalence (associée à une partition).}{\ChapRef{chap:axiomatique}}
\Not{\preceq}{ordre (partiel) de raffinement sur les partitions.}{\ChapRef{chap:axiomatique}}
\Not{u}{ultramétrique.}{\ChapRef{chap:axiomatique}}
\Not{\theta}{dendrogramme ou clustering hiérarchique.}{\ChapRef{chap:axiomatique}}
\Not{\theta(r)}{partition ou recouvrement produit au niveau \(r\).}{\ChapRef{chap:axiomatique}}
\Not{\theta^{\mathrm{SL}}}{hiérarchie produite par le Single-Linkage.}{\ChapRef{chap:single-linkage_trois_vues}}
\Not{\theta^{\mathrm{HGP}}_K}{hiérarchie produite par \HGP-Clusterer à l'ordre \(K\).}{\ChapRef{chap:hgp_clusterer}}
\end{NotationTable}

\begin{NotationTable}{Distances de liaison}
\Not{C,C'}{clusters, amas ou composantes connexes.}{\CtxGlobal}
\Not{C^{\mathrm{discret}}}{cluster discret associé à une composante continue de forte densité.}{\ChapRefs{chap:single-linkage_trois_vues}{chap:hgp_clusterer}}
\Not{d_{\mathrm{SL}}}{distance de Single-Linkage.}{\ChapRef{chap:single-linkage_trois_vues}}
\Not{d_{\mathrm{mreach},K}}{distance d'atteignabilité mutuelle du Robust Single-Linkage.}{\ChapRef{chap:limites_single_linkage}}
\Not{r_t}{rayon de fusion à l'étape \(t\) du Single-Linkage.}{\ChapRef{chap:single-linkage_trois_vues}}
\Not{r_{\mathrm{connect}}}{rayon auquel tous les points sont réunis dans une seule composante.}{\ChapRef{chap:single-linkage_trois_vues}}
\Not{\mathrm{loss}(C)}{fonction de coût utilisée pour explorer une hiérarchie.}{\ChapRef{chap:hacker_hdbscan}}
\Not{\widetilde E(C)}{stabilité ou excès de masse relatif d'un cluster.}{\ChapRef{chap:limites_single_linkage}}
\Not{\omega : x \rightsquigarrow x'}{Chemin $\omega$ au sein d'un graphe (liste d'arêtes) menant du sommet $x$ au sommet $x'$.}{\CtxGlobal}
\end{NotationTable}

\begin{NotationTable}{Graphes et arbres couvrants}
\Not{G=(V,E)}{graphe de sommets \(V\) et d'arêtes \(E\).}{\CtxGlobal}
\Not{G=(V,E,w)}{graphe pondéré par une fonction de poids \(w\).}{\ChapRef{chap:single-linkage_trois_vues}}
\Not{w(e)}{poids de l'arête \(e \in E\).}{\ChapRef{chap:single-linkage_trois_vues}}
\Not{\mathrm{MST}}{arbre minimum couvrant.}{\ChapRefsThree{chap:single-linkage_trois_vues}{chap:generalisation_mst}{chap:frangi_fissures}}
\Not{\mathrm{MST}_{\le r}}{arbre minimum couvrant élagué au niveau \(r\).}{\ChapRef{chap:single-linkage_trois_vues}}
\Not{G_{\le r}}{sous-graphe contenant les arêtes de poids au plus \(r\).}{\ChapRef{chap:single-linkage_trois_vues}}
\Not{\mathcal G(\mathcal X,r)}{graphe géométrique de rayon \(r\) sur \(\mathcal X\).}{\ChapRefs{chap:single-linkage_trois_vues}{chap:hgp_clusterer}}
\Not{\mathrm{RNG}}{graphe de voisinage relatif.}{\ChapRef{chap:single-linkage_trois_vues}}
\end{NotationTable}

\begin{NotationTable}{Densité et estimateur \(K\)-NN}
\Not{K}{nombre de voisins pour l'estimateur \KNN.}{\ChapRefs{chap:single-linkage_trois_vues}{chap:hgp_clusterer}}
\Not{f}{densité théorique.}{\ChapRef{chap:single-linkage_trois_vues}}
\Not{\widehat f_K}{estimateur de densité aux \(K\) plus proches voisins.}{\ChapRefs{chap:single-linkage_trois_vues}{chap:hgp_clusterer}}
\Not{\widehat\lambda}{notation abrégée pour \(\widehat f_K(y)\).}{\ChapRef{chap:single-linkage_trois_vues}}
\Not{\lambda}{niveau de densité, ou intensité d'un processus de \nom{Poisson}.}{\ChapRefs{chap:single-linkage_trois_vues}{sec:chap_percolation}}
\Not{r}{rayon associé au niveau de densité par changement de variable.}{\CtxIntro}
\Not{r_K(y)}{distance de \(y\) à son \(K\)-ième plus proche voisin dans \(\mathcal X\).}{\ChapRef{chap:single-linkage_trois_vues}}
\Not{L_f(\lambda)}{ensemble de niveau supérieur \(\{x\mid f(x)\ge \lambda\}\).}{\ChapRef{chap:single-linkage_trois_vues}}
\Not{L_K(r)}{ensemble de niveau supérieur associé à \(\widehat f_K\), paramétré par le rayon.}{\ChapRefs{chap:single-linkage_trois_vues}{chap:hgp_clusterer}}
\Not{H_f(\lambda)}{amas de forte densité de \(f\), c'est-à-dire \(\pi_0(L_f(\lambda))\).}{\ChapRef{chap:single-linkage_trois_vues}}
\Not{H_{\widehat f_K}^{\mathrm{discrets}}(r)}{amas discrets de forte densité associés à l'estimateur \KNN.}{\ChapRefs{chap:single-linkage_trois_vues}{chap:hgp_clusterer}}
\end{NotationTable}

\begin{NotationTable}{Complexes simpliciaux et \HGP-Clusterer}
\Not{K}{ordre des interactions dans \HGP-Clusterer.}{\ChapRef{chap:hgp_clusterer}}
\Not{\sigma,\tau}{simplexes, vus comme sous-ensembles finis de \(\mathcal X\).}{\ChapRefs{chap:hgp_clusterer}{chap:generalisation_mst}}
\Not{\check C(\mathcal X,r)}{complexe de \nom[Cech]{Čech} de rayon \(r\) sur \(\mathcal X\).}{\ChapRef{chap:hgp_clusterer}}
\Not{\check C_{K-1}(\mathcal X,r)}{ensemble des \((K-1)\)-simplexes de \(\check C(\mathcal X,r)\).}{\ChapRef{chap:hgp_clusterer}}
\Not{T_r(\sigma)}{région témoin du simplexe \(\sigma\), égale à \(\bigcap_{x\in\sigma}\overline B(x,r)\).}{\ChapRef{chap:hgp_clusterer}}
\Not{\Gamma_K(\mathcal X,r)}{graphe auxiliaire des \((K-1)\)-simplexes.}{\ChapRef{chap:hgp_clusterer}}
\Not{\Gamma_K(\mathcal X)_{\le r}}{version élaguée de \(\Gamma_K\) au niveau \(r\).}{\ChapRef{chap:generalisation_mst}}
\Not{\rho(\sigma)}{rayon de naissance du simplexe \(\sigma\).}{\ChapRef{chap:generalisation_mst}}
\Not{B_\sigma=\overline B(c_\sigma,\rho(\sigma))}{plus petite boule fermée englobante de \(\sigma\).}{\ChapRef{chap:generalisation_mst}}
\Not{S(A)}{ensemble de support de la plus petite boule englobante de \(A\).}{\ChapRef{chap:generalisation_mst}}
\Not{\mathcal T_K}{\(K\)-arbre minimum couvrant.}{\ChapRef{chap:generalisation_mst}}
\end{NotationTable}

\begin{NotationTable}{Géométrie d'ordre supérieur et sortie en partition}
\Not{\mathrm{Del}_K(\mathcal X)}{mosaïque de \nom{Delaunay} d'ordre \(K\).}{\ChapRef{chap:generalisation_mst}}
\Not{\mathcal F_K}{ensemble des \((K-1)\)-simplexes construits par l'algorithme.}{\ChapRef{chap:conclusion_partie_II}}
\Not{\ell(\tau)}{étiquette de cluster portée par la face \(\tau\).}{\ChapRef{chap:conclusion_partie_II}}
%\Not{q_{\mathrm{clust}}}{nombre de clusters sélectionnés.}{\ChapRef{chap:conclusion_partie_II}}
%\Not{S_\tau}{score local positif associé à une face \(\tau\).}{\ChapRef{chap:conclusion_partie_II}}
\Not{\psi(t)}{fonction de poids décroissante, souvent \(\psi(t)=1/t^p\).}{\ChapRef{chap:conclusion_partie_II}}
%\Not{V_x(c)}{score de vote du point \(x\) pour le cluster \(c\).}{\ChapRef{chap:conclusion_partie_II}}
\Not{\widehat\ell(x)}{étiquette finale attribuée au point \(x\).}{\ChapRef{chap:conclusion_partie_II}}
\Not{\widehat C_c}{ensemble des points finalement étiquetés par \(c\).}{\ChapRef{chap:conclusion_partie_II}}
\end{NotationTable}

\begin{NotationTable}{Complexe de \nom[Vietoris-Rips]{Vietoris--Rips} et UMAP}
\Not{\mathrm{VR}(\mathcal X,r)}{complexe de \nom[Vietoris-Rips]{Vietoris--Rips} de rayon \(r\).}{\ChapRef{chap:conclusion_partie_II}}
%\Not{G_r=(\mathcal X,E_r)}{graphe seuil associé à \(\mathrm{VR}(\mathcal X,r)\).}{\ChapRef{chap:conclusion_partie_II}}
\Not{E_r}{arêtes \(\{x,x'\}\) telles que \(d(x,x')\le 2r\).}{\ChapRef{chap:conclusion_partie_II}}
\Not{\alpha_p}{constante d'intercalage %\nom[Cech]{Čech}--\nom[Vietoris-Rips]{Vietoris--Rips}, 
\(\sqrt{2p/(p+1)}\).}{\ChapRef{chap:conclusion_partie_II}}
\Not{k_{\mathrm{nn}}}{nombre de voisins dans UMAP.}{\ChapRef{chap:conclusion_partie_II}}
\Not{\rho_i,\sigma_i}{distance locale minimale et échelle locale du point \(x_i\) dans UMAP.}{\ChapRef{chap:conclusion_partie_II}}
\Not{p_{j|i}}{poids orienté de voisinage de \(x_j\) vu depuis \(x_i\).}{\ChapRef{chap:conclusion_partie_II}}
\Not{w_{ij}}{poids symétrisé par union floue dans UMAP.}{\ChapRef{chap:conclusion_partie_II}}
\Not{\mathcal L_{\mathrm{UMAP}}(Y)}{entropie croisée optimisée par UMAP.}{\ChapRef{chap:conclusion_partie_II}}
\end{NotationTable}

\begin{NotationTable}{Percolation continue}
\Not{\mathcal X_\lambda}{processus ponctuel de \nom{Poisson} homogène d'intensité \(\lambda\).}{\ChapRef{sec:chap_percolation}}
\Not{\mathcal X_\lambda^0}{version de \nom{Palm}, \(\mathcal X_\lambda\cup\{0\}\).}{\ChapRef{sec:chap_percolation}}
\Not{\xi}{configuration ponctuelle localement finie.}{\ChapRef{sec:chap_percolation}}
\Not{\mu}{intensité normalisée, \(\mu=\lambda\omega_p2^{-p}\).}{\ChapRef{sec:chap_percolation}}
\Not{\mathrm{cc}}{famille de composantes connexes : \(\mathrm{poly}\), \(\mathrm{core}\) ou \(\mathrm{dbscan}\).}{\ChapRef{sec:chap_percolation}}
\Not{\mathrm{poly}}{\(K\)-polyèdres.}{\ChapRef{sec:chap_percolation}}
\Not{\mathrm{core}}{composantes de cœurs du Robust Single-Linkage.}{\ChapRef{sec:chap_percolation}}
\Not{\mathrm{dbscan}}{composantes de type \DBSCAN.}{\ChapRef{sec:chap_percolation}}
\Not{\Theta^{\mathrm{cc}}_{K,p,r}(\lambda)}{probabilité de percolation de la méthode \(\mathrm{cc}\).}{\ChapRef{sec:chap_percolation}}
\Not{\widehat\Theta^{\mathrm{cc}}}{estimateur empirique de la fonction de percolation.}{\ChapRef{sec:chap_percolation}}
\end{NotationTable}

\begin{NotationTable}{Seuils et vitesse de percolation}
%\Not{C_\infty}{composante infinie ou composante géante.}{\ChapRef{sec:chap_percolation}}
\Not{\mathcal C_\infty(A)}{réunion des composantes connexes non bornées de \(A\).}{\ChapRef{sec:chap_percolation}}
\Not{Q=[0,1]^p}{hypercube unité.}{\ChapRef{sec:chap_percolation}}
\Not{\lambda_c^{\mathrm{cc}}(K,p)}{seuil critique de percolation.}{\ChapRef{sec:chap_percolation}}
\Not{\lambda_\alpha^{\mathrm{cc}}(K,p)}{quantile de niveau \(\alpha\) de la courbe de percolation.}{\ChapRef{sec:chap_percolation}}
\Not{\lambda^{\mathrm{cc}}_{1-\varepsilon}(K,p)}{quantile nécessaire pour atteindre un rappel \(1-\varepsilon\).}{\ChapRef{sec:chap_percolation}}
\Not{v^{\mathrm{cc}}_{\mathrm{crit},K,p}(\varepsilon)}{vitesse critique de percolation \(\lambda_c/\lambda_{1-\varepsilon}\).}{\ChapRef{sec:chap_percolation}}
\Not{v^{\mathrm{cc}}_{K,p}(\varepsilon)}{vitesse empirique \(\lambda_\varepsilon/\lambda_{1-\varepsilon}\).}{\ChapRef{sec:chap_percolation}}
\Not{\varepsilon}{niveau de queue, associé au haut rappel \(1-\varepsilon\).}{\ChapRef{sec:chap_percolation}}
%\Not{\alpha}{niveau de quantile d'une fonction de percolation.}{\ChapRef{sec:chap_percolation}}
\end{NotationTable}

\begin{NotationTable}{Fenêtre gaussienne \(K\to+\infty\)}
\Not{N_\mu(x)}{nombre de points dans la boule de rayon \(1/2\) autour de \(x\).}{\ChapRef{sec:chap_percolation}}
\Not{Z_\mu(x)}{champ de comptage centré réduit, \((N_\mu(x)-\mu)/\sqrt\mu\).}{\ChapRef{sec:chap_percolation}}
\Not{G_p}{champ gaussien stationnaire limite en dimension \(p\).}{\ChapRef{sec:chap_percolation}}
\Not{\rho_p(z)}{covariance du champ \(G_p\).}{\ChapRef{sec:chap_percolation}}
\Not{a}{paramètre de fluctuation dans la fenêtre \(\mu=K+a\sqrt K\).}{\ChapRef{sec:chap_percolation}}
\Not{\mu_K(a)}{intensité normalisée \(K+a\sqrt K\).}{\ChapRef{sec:chap_percolation}}
\Not{\lambda_K(a)}{intensité non normalisée associée à \(\mu_K(a)\).}{\ChapRef{sec:chap_percolation}}
%\Not{F_K^{\mathrm{cc}}(a)}{courbe de percolation réparamétrée par \(a\).}{\ChapRef{sec:chap_percolation}}
\Not{E_a}{excursion gaussienne \(\{x\in\mathbb R^p\mid G_p(x)\ge -a\}\).}{\ChapRef{sec:chap_percolation}}
\Not{a_c^{\mathrm{cc}}}{seuil critique dans l'échelle gaussienne \(a\).}{\ChapRef{sec:chap_percolation}}
%\Not{\mathcal H_p}{espace de \nom{Cameron--Martin} ou RKHS associé à \(G_p\).}{\ChapRef{sec:chap_percolation}}
\end{NotationTable}

\begin{NotationTable}{Cadre bayésien pour les communautés}
\Not{K}{nombre de communautés.}{\ChapRef{chap:bayes_clusters_graphes}}
\Not{V_n}{ensemble des sommets, \(V_n=\{1,\ldots,n\}\).}{\ChapRef{chap:bayes_clusters_graphes}}
\Not{E_n}{ensemble des arêtes potentielles.}{\ChapRef{chap:bayes_clusters_graphes}}
\Not{\Lab_K}{ensemble des étiquettes de communautés, \(\{1,\ldots,K\}\).}{\ChapRef{chap:bayes_clusters_graphes}}
\Not{\Conf_n}{espace des configurations, \((\Lab_K)^{V_n}\).}{\ChapRef{chap:bayes_clusters_graphes}}
\Not{\sigma,\tau}{configurations de communautés (fonctions de $V_n$ vers $\Lab_K$).}{\ChapRef{chap:bayes_clusters_graphes}}
\Not{\sigma_i}{étiquette du sommet \(i\) dans la configuration \(\sigma\).}{\ChapRef{chap:bayes_clusters_graphes}}
\Not{\Sigma_n}{configuration cachée à retrouver.}{\ChapRef{chap:bayes_clusters_graphes}}
\Not{\Wspace}{espace des poids d'arêtes.}{\ChapRef{chap:bayes_clusters_graphes}}
\Not{X_n}{attributs observés sur les sommets.}{\ChapRef{chap:bayes_clusters_graphes}}
\Not{W_n}{poids observés sur les arêtes.}{\ChapRef{chap:bayes_clusters_graphes}}
\end{NotationTable}

\begin{NotationTable}{Postérieure et recouvrement}
\Not{\mu_{0,n}}{loi \emph{a priori} sur les configurations.}{\ChapRef{chap:bayes_clusters_graphes}}
\Not{Q_n}{canal d'observation des poids d'arêtes.}{\ChapRef{chap:bayes_clusters_graphes}}
\Not{q_n(w\mid x,\sigma)}{densité du canal \(Q_n\).}{\ChapRef{chap:bayes_clusters_graphes}}
\Not{\Obs_n=(X_n,W_n)}{observation disponible.}{\ChapRef{chap:bayes_clusters_graphes}}
\Not{\mu_{n,x,w}}{loi postérieure de \(\Sigma_n\) sachant \(X_n=x,W_n=w\).}{\ChapRef{chap:bayes_clusters_graphes}}
\Not{Z_n(x,w)}{constante de normalisation de la postérieure.}{\ChapRef{chap:bayes_clusters_graphes}}
\Not{\tau_n}{sortie d'un algorithme de reconstruction.}{\ChapRef{chap:bayes_clusters_graphes}}
\Not{\rho_n}{loi d'un tirage au hasard, indépendante de l'observation.}{\ChapRef{chap:bayes_clusters_graphes}}
\Not{\Phi_n}{score comparant vérité cachée et sortie d'algorithme.}{\ChapRef{chap:bayes_clusters_graphes}}
\Not{\operatorname{ov}_n(\sigma,\tau)}{recouvrement (\emph{overlap}) maximal entre deux configurations, modulo permutation.}{\ChapRef{chap:bayes_clusters_graphes}}
\Not{\RG_n(s)}{probabilité de succès du meilleur tirage au hasard (\emph{random guess}) au seuil \(s\).}{\ChapRef{chap:bayes_clusters_graphes}}
\end{NotationTable}

\begin{NotationTable}{Mesures de \nom{Gibbs} et dynamiques de clusters}
\Not{U_{x,w}(\sigma)}{énergie postérieure associée à \((x,w)\).}{\ChapRef{chap:bayes_clusters_graphes}}
\Not{\mu(\sigma)}{mesure de \nom{Gibbs} sur les configurations.}{\ChapRef{chap:bayes_clusters_graphes}}
\Not{Z}{fonction de partition d'une mesure de \nom{Gibbs}.}{\ChapRef{chap:bayes_clusters_graphes}}
\Not{\mathcal B}{famille des liens ou interactions locales.}{\ChapRef{chap:bayes_clusters_graphes}}
\Not{b}{lien local (\emph{bond}) : arête, triangle, tétraèdre ou motif plus général.}{\ChapRef{chap:bayes_clusters_graphes}}
\Not{U_b(\sigma)}{contribution énergétique du lien \(b\).}{\ChapRef{chap:bayes_clusters_graphes}}
\Not{\omega_b}{sous-lien gelé tiré sur le lien \(b\).}{\ChapRef{chap:bayes_clusters_graphes}}
\Not{P_b^\sigma}{loi de gel du lien \(b\) sous la configuration \(\sigma\).}{\ChapRef{chap:bayes_clusters_graphes}}
\Not{T(\sigma,\sigma')}{probabilité de transition de la chaîne de \nom{Markov} de type \nom{Swendsen--Wang} généralisée.}{\ChapRef{chap:bayes_clusters_graphes}}
\Not{\kappa_n}{objet gelé produit par un pas de dynamique de clusters.}{\ChapRef{chap:bayes_clusters_graphes}}
\Not{\mathcal C(\kappa_n)}{clusters induits par les gels de \(\kappa_n\).}{\ChapRef{chap:bayes_clusters_graphes}}
\Not{\theta_\delta(\kappa_n)}{masse portée par les clusters de taille au moins \(\delta n\).}{\ChapRef{chap:bayes_clusters_graphes}}
\end{NotationTable}

\begin{NotationTable}{Modèle de \nom{Sankararaman--Baccelli}}
\Not{H}{sous-graphe observé.}{\ChapRef{chap:bayes_clusters_graphes}}
\Not{f_{\mathrm{in}},f_{\mathrm{out}}}{fonctions de connexion intra-communauté et inter-communautés.}{\ChapRef{chap:bayes_clusters_graphes}}
\Not{r_e}{longueur de l'arête potentielle \(e\).}{\ChapRef{chap:bayes_clusters_graphes}}
\Not{W_e(H)}{poids signé déduit de l'observation \(H\).}{\ChapRef{chap:bayes_clusters_graphes}}
\Not{\mu_H}{postérieure de \nom{Gibbs} sachant le sous-graphe observé \(H\).}{\ChapRef{chap:bayes_clusters_graphes}}
\Not{p}{paramètre homogène de connexion sur la grille triangulaire.}{\ChapRef{chap:bayes_clusters_graphes}}
\Not{u_p}{poids absolu \(\log(p/(1-p))\).}{\ChapRef{chap:bayes_clusters_graphes}}
\Not{p_{\mathrm{lien}}}{probabilité effective de conservation d'une arête gelée.}{\ChapRef{chap:bayes_clusters_graphes}}
\Not{p_c^{\mathrm{edge}}}{seuil critique obtenu avec la dynamique arête par arête.}{\ChapRef{chap:bayes_clusters_graphes}}
\Not{p_c^\triangle}{seuil critique obtenu avec la dynamique triangulaire.}{\ChapRef{chap:bayes_clusters_graphes}}
\end{NotationTable}

\begin{NotationTable}{Images et filtre de \nom{Frangi}}
\Not{\Omega}{domaine discret de l'image, typiquement \(\Omega\subset\mathbb Z^2\).}{\ChapRef{chap:frangi_fissures}}
\Not{I:\Omega\to\mathbb R}{image scalaire.}{\ChapRef{chap:frangi_fissures}}
\Not{\mathcal M}{ensemble des modalités.}{\ChapRef{chap:frangi_fissures}}
\Not{I^{(m)}}{image de la modalité \(m \in \mathcal M\).}{\ChapRef{chap:frangi_fissures}}
\Not{\Sigma}{ensemble des échelles gaussiennes du filtre de \nom{Frangi}.}{\ChapRef{chap:frangi_fissures}}
\Not{G_\sigma}{noyau gaussien d'échelle \(\sigma\).}{\ChapRef{chap:frangi_fissures}}
\Not{\Hess_\sigma I(x)}{hessienne lissée de l'image \(I\) à l'échelle \(\sigma\).}{\ChapRef{chap:frangi_fissures}}
\Not{\lambda_1^\sigma(x),\lambda_2^\sigma(x)}{valeurs propres ordonnées de la hessienne lissée.}{\ChapRef{chap:frangi_fissures}}
%\Not{V_\sigma(x)}{réponse de \nom{Frangi} à l'échelle \(\sigma\).}{\ChapRef{chap:frangi_fissures}}
%\Not{V(x)}{réponse multi-échelle maximale de \nom{Frangi}.}{\ChapRef{chap:frangi_fissures}}
%\Not{\beta,c}{paramètres d'échelle de la réponse de \nom{Frangi}.}{\ChapRef{chap:frangi_fissures}}
\end{NotationTable}

\begin{NotationTable}{Graphe de \nom{Frangi} et squelette}
\Not{\widehat\Gamma}{squelette ou réseau de fissures extrait.}{\ChapRef{chap:frangi_fissures}}
\Not{V\subseteq\Omega}{ensemble des pixels candidats.}{\ChapRef{chap:frangi_fissures}}
\Not{R}{rayon spatial de voisinage entre pixels.}{\ChapRef{chap:frangi_fissures}}
\Not{E_R}{arêtes reliant les pixels candidats à distance au plus \(R\).}{\ChapRef{chap:frangi_fissures}}
\Not{\theta^\sigma(x)}{orientation locale estimée à partir de la hessienne.}{\ChapRef{chap:frangi_fissures}}
\Not{S_{ij}}{similarité de \nom{Frangi} multi-échelle entre deux pixels \(x_i\) et \(x_j\).}{\ChapRef{chap:frangi_fissures}}
\Not{d_{ij}}{dissimilarité (de \nom{Frangi}) géométrique associée à \(S_{ij}\).}{\ChapRef{chap:frangi_fissures}}
\Not{\mathcal T}{arbre couvrant minimum utilisé pour extraire le squelette.}{\ChapRef{chap:frangi_fissures}}
\Not{\tau}{seuil relatif de conservation des pixels ou arêtes.}{\ChapRef{chap:frangi_fissures}}
\Not{\tau_c}{seuil relatif de centralité.}{\ChapRef{chap:frangi_fissures}}
\Not{\Hess_\sigma^{\mathrm{fused}}(x)}{hessienne fusionnée au pixel \(x\).}{\ChapRef{chap:frangi_fissures}}
%\Not{G^\triangle}{graphe dual des arêtes pour les composantes triangles-connexes.}{\ChapRef{chap:frangi_fissures}}
\end{NotationTable}

\begin{NotationTable}{Métriques d'évaluation}
\Not{A}{vérité terrain vue comme ensemble de pixels.}{\ChapRef{chap:frangi_fissures}}
\Not{B}{prédiction vue comme ensemble de pixels.}{\ChapRef{chap:frangi_fissures}}
\Not{J(A,B)}{indice de \nom{Jaccard}, \(|A\cap B|/|A\cup B|\).}{\ChapRef{chap:frangi_fissures}}
\Not{T(A,B)}{indice de \nom{Tversky}.}{\ChapRef{chap:frangi_fissures}}
\Not{\alpha,\beta}{pondérations de l'indice de \nom{Tversky}.}{\ChapRef{chap:frangi_fissures}}
%\Not{W_1(A,B)}{distance de \nom{Wasserstein} entre distributions discrètes portées par deux réseaux.}{\ChapRef{chap:frangi_fissures}}
\Not{\mathrm{IoU}}{autre nom de l'indice de \nom{Jaccard} (\emph{Intersection over Union}).}{\ChapRef{chap:frangi_fissures}}
\Not{\mathrm{ARI}}{indice de \nom{Rand} ajusté (\emph{Adjusted \nom{Rand} Index}).}{\ChapRef{chap:conclusion_partie_II}}
\end{NotationTable}

\endgroup